\section{Formal Framework}
\label{sec:definitions}

\subsection{Admissible histories}

Let $\mathcal{H}$ denote the set of admissible histories of a physical system. Elements
$h \in \mathcal{H}$ represent complete realizations of the system subject to underlying
physical constraints. No preferred global temporal parameterization is assumed on
$\mathcal{H}$; however, histories may support partial orderings induced by experimental
interventions or boundary conditions.

Throughout, $\Sigma_1$ (``Sigma-one'') denotes algebraically closed reduced descriptions, while
$\Sigma_2$ (``Sigma-two'') denotes trajectory-level admissibility structures.

\subsection{Operational reduction}

An operational description is obtained via a reduction map
\begin{equation}
\Pi : \mathcal{H} \rightarrow \mathcal{R},
\end{equation}
where $\mathcal{R}$ is a space of reduced descriptions, such as density matrices, quantum
channels, or static multi-time process tensors. The reduction $\Pi$ is generally many-to-one
and represents the information retained by a given experimental or theoretical interface.

A reduction is said to be \emph{degenerate} if there exist distinct histories
$h_1 \neq h_2 \in \mathcal{H}$ such that
\begin{equation}
\Pi(h_1) = \Pi(h_2).
\end{equation}

\subsection{Algebraically closed reduced descriptions (\texorpdfstring{$\Sigma_1$}{Sigma-1})}

A reduced description $r \in \mathcal{R}$ is said to be \emph{algebraically closed} if all
admissible future predictions are determined by a fixed algebraic object together with
composition rules internal to a prescribed algebraic structure. Examples include
instantaneous quantum states with completely positive maps, quantum channels, and static
multi-time process tensors equipped with contraction rules.

We refer to any framework relying exclusively on algebraically closed reduced descriptions
as a $\Sigma_1$ framework.

\subsection{Interventions and history composition}

Let $\mathcal{I}$ denote a specified class of physical interventions. Interventions are
externally applied operations performed after the reduction, acting only on degrees of
freedom represented in $\mathcal{R}$. Examples include bounded Hamiltonian control pulses,
completely positive maps, and standard measurement-and-feedback operations. Interventions
that directly access discarded degrees of freedom are excluded.

We assume that admissible histories support a composition relation with interventions.
For $h,h' \in \mathcal{H}$ and $I \in \mathcal{I}$, we write
\begin{equation}
h \sim_I h'
\end{equation}
to mean that $h'$ is a possible history resulting from applying intervention $I$ to history
$h$. No determinism or uniqueness is assumed; $\sim_I$ is a relation, not a function. 

Measurements performed on the retained system $S$ that do not selectively post-condition on
environment outcomes are included as admissible interventions; operations requiring access to,
or conditioning on, the state of discarded degrees of freedom are excluded.

We do not require the intervention class $\mathcal{I}$ to be closed under composition or to form
a group. However, for multi-step counterfactual reasoning, one would typically assume that if
$I_1, I_2 \in \mathcal{I}$, then their sequential composition $I_2 \circ I_1$ is also admissible.
All results of this work apply whether or not such closure holds.

\subsection{Continuation sets}

For a reduced description $r \in \mathcal{R}$, define the \emph{continuation set} under
$\mathcal{I}$ as
\begin{equation}
\mathrm{Cont}(r;\mathcal{I})
=
\left\{
(I, r') \;\middle|\;
\exists\, h,h' \in \mathcal{H} \text{ such that }
\Pi(h)=r,\; h \sim_I h',\; \Pi(h')=r'
\right\}.
\end{equation}

The continuation set therefore records, for each admissible intervention, the set of
possible reduced outcomes compatible with at least one underlying history consistent
with $r$. This definition makes explicit that continuation structure is inherited from
$\mathcal{H}$ via the reduction $\Pi$, without assuming any reduced-level dynamics.

\subsection{Faithfulness with respect to interventions}

\begin{definition}[Faithfulness]
A reduction $\Pi$ is \emph{faithful with respect to} $\mathcal{I}$ if, for all histories
$h_1,h_2 \in \mathcal{H}$,
\begin{equation}
\Pi(h_1) = \Pi(h_2)
\;\Rightarrow\;
\mathrm{Cont}(\Pi(h_1);\mathcal{I})
=
\mathrm{Cont}(\Pi(h_2);\mathcal{I}).
\end{equation}
\end{definition}

Faithfulness requires that identical reduced descriptions support identical sets of
counterfactual continuations under the specified intervention class.

\subsection{Counterfactual prediction (operational sense)}

We say that a reduced description $r \in \mathcal{R}$ supports \emph{counterfactual prediction}
for an intervention class $\mathcal{I}$ if, for each $I \in \mathcal{I}$, the reduced theory
assigns a unique (or uniquely specified) set of admissible reduced outcomes
\begin{equation}
\widehat{\mathrm{Out}}(r;I) \subseteq \mathcal{R}
\end{equation}
that is invariant over the preimage $\Pi^{-1}(r)$; equivalently, if $\mathrm{Cont}(r;\mathcal{I})$
is well-defined independent of which history in $\Pi^{-1}(r)$ is realized. In an unfaithful cut,
identical reduced descriptions correspond to different counterfactual responses to the same $I$,
so no Σ$_1$ dynamics defined solely on $\mathcal{R}$ can determine the correct answer.

\subsection{The unfaithful cut}

A reduction $\Pi$ constitutes an \emph{unfaithful cut} if it is degenerate and fails to be
faithful with respect to $\mathcal{I}$. In this case, the reduced description is structurally
incapable of supporting counterfactual prediction under $\mathcal{I}$, regardless of the
dynamical laws imposed at the reduced level.